\subsection{Two binary stalks and five intermediate arrows}

\begin{proposition}[Proved: five labels cannot support a source cycle]
\label{A7:two-binary-five-labels}
Let $\mathcal C$ be a finite acyclic category with full support
$0<1<2<3<4$, and write $h_{ij}=|\mathcal C(i,j)|$.
Every rational degree-two cohomology product of $B\mathcal C$
is zero if
\[
 h_{01}=h_{02}=2,\qquad h_{13}\leq5.
\]
The same conclusion holds under the dual conditions
$h_{34}=h_{24}=2$ and $h_{13}\leq5$.
There is no restriction on the two long hom-sets or on $h_{04}$.
\end{proposition}
\begin{proof}
Use $A=\mathcal C(0,1)$, $U=\mathcal C(0,2)$,
$B=\mathcal C(1,2)$, $V=\mathcal C(1,3)$ and
$C=\mathcal C(2,3)$, and let
$P:A\times B\to U$, $Q:B\times C\to V$ be composition.
If every map $P(-,b):A\to U$ is constant, the existing
short-composite obstruction, Proposition~\ref{A7:short-composite-gates},
proves the assertion. Otherwise some $b_*\in B$ induces a
bijection, since both $A$ and $U$ have two elements.

Work first in the universal short-product category $F(P,Q,R)$.
Its $03$ hom-set $L$ is the quotient of
$(A\times V)\amalg(U\times C)$ by the relations
\[
 (a,Q(b,c))\sim(P(a,b),c).
\]
On the vertex set $L$, form the simple undirected graph $G$
whose edges are the distinct nonconstant image pairs of all
maps $A\to L$ indexed by $V$ and all maps $U\to L$ indexed
by $C$. Loops are discarded and parallel image pairs give one
edge. We show that a cycle in $G$ would require at least six
\emph{distinct} elements of $V$.

To keep track of the quotient, initially take one abstract
interval $I_v$ with endpoint set $A\times\{v\}$ for each
$v\in V$, and one interval $J_c$ with endpoint set
$U\times\{c\}$ for each $c\in C$.
First impose only the relations belonging to bijective columns
$P(-,b)$. Each such relation identifies both endpoints of
$J_c$ with the corresponding endpoints of $I_{Q(b,c)}$,
either preserving or reversing their order. Partition all these
intervals into classes under these whole-interval identifications.
Within each class the endpoint quotient has either two vertices
or one vertex. The latter case includes an orientation-reversing
identification around a chain of intervals. Regard a two-vertex
class as a single interval and a one-vertex class as a point;
no topological quotient with an extra midpoint is being used.
Different classes at this stage have disjoint vertex sets.
Every class contains an interval $I_v$, because the chosen
bijective column $b_*$ attaches every $J_c$ to such an interval.

The remaining columns are constant. If $P(-,b)$ has value
$u\in U$, their relations identify both endpoints of
$I_{Q(b,c)}$ with the $u$ endpoint of $J_c$. Mark every
interval class which contains an $I_v$ occurring on this
constant-column side. Each marked class is forced to collapse
to a single final vertex. After these collapses, every further
identification between different initial classes is incident to
the point of a marked class. All equivalences used here are
exactly the displayed defining endpoint relations for $L$.
In particular, the vertices of $G$ are their final vertex classes.

A nonloop edge of $G$ is the image of an uncollapsed interval
class. It therefore contains a label $v\in V$ belonging to
that class, and different edges can be represented by different
classes and hence different labels. A marked class cannot
represent a nonloop edge.
If a final vertex is incident to edges coming from two different
interval classes, its endpoint equivalence class must contain
a marked class: before the constant-column relations the two
classes were disjoint, and every relation joining different
classes is incident to a marked class. A marked class has just
one final vertex, even if it anchors arbitrarily many merges.
Consequently different such final vertices require different
marked classes and different labels from $V$.

Suppose now that $G$ has a simple cycle of length $m\geq3$.
Its $m$ distinct edges require $m$ distinct labels from
uncollapsed classes. At each of its $m$ distinct vertices,
the two incident cycle edges come from different interval
classes, so that vertex requires a label from a marked class.
Those $m$ marked labels are pairwise distinct and are disjoint
from the $m$ edge labels. Thus $|V|\geq2m\geq6$, a
contradiction. Reversed identifications were accounted for by
the possible point classes; parallel edges cannot create a
cycle in the simple graph relevant to the tree condition.
This proves that $G$ is a forest, including all its isolated
vertices.

Extend this forest to a tree on the full finite set $L$.
Use the unique two-vertex trees on $A$ and $U$, and the point
on the singleton extreme stalk $F(0,4)$. The maps $A\to U$
carry an edge to an edge or a vertex. The same holds for every
map from $A$ or $U$ to $L$ by the definition of $G$, and all
maps to the final singleton are constant. These are all the
source-act maps, so Lemma~\ref{A7:tree-act} proves that all
positive rational products on $BF$ vanish. The graph argument
did not use $R$. The universal transfer theorem,
Theorem~\ref{A7:universal-short-all-fields}, therefore proves
the desired degree-two product vanishing for every completion
and hence for $\mathcal C$. Reversing the object order gives
the stated dual condition.
\end{proof}
